\scanpage{239}
\mspage{122}

3. Conditions d'homogénéité.

\deleted{Nous nous intéressons maintenant au cas où} Revenons au cas de \deleted{$G_0$ a la propriété d'homogénéité} \deleted{\unsure{$\Gamma$}}

% slanted
\margin{Cas d'un}

\[
M=(\Gamma,\quad Z,\quad\pi),
\]

je dis qu'on a équivalence entre les propriétés

\[
\left\{\begin{aligned}
&\text{a) }G_0\text{ est transitif sur }R\\
&\text{b) }G_0\text{ est transitif sur }R_{\varphi_0}\ \text{i.e. sur l'ens }\overrightarrow{A_{\varphi_0}}\\
&\qquad\text{des arêtes orientées du graphe }(S,A_{\varphi_0}),\ \text{pour}\\
&\qquad\text{\uline{un} }\varphi_0\in\pi\\
&\text{b$'$) Comme b), pour tout }\varphi\in\pi
\end{aligned}\right.
\tag*{(59)}
\]

% slanted
\margin{\uline{NB} a) s'énonce sans référence à un $\pi$ contrairement à b) b$'$) --- c'est une propriété d'un \deleted{\unsure{$\pi$ mais}} groupe d'autom $G_0$ d'un prégraphe, en \deleted{fait} \inserted{\unsure{présence|défini}} d'un $Z$. (Mais on suppose ici $G_0$ en termes d'un $Z,\pi$).}

L'équivalence des a) et b) \inserted{ou a) et b$'$)} résulte des \deleted{\unsure{ceci} fait}

\[
R\simeq R_{\varphi_0}\qquad(G_0\text{-isom})\quad\text{via }\mathrm{id}\times\varphi_0\,\unsure{.}
\]

\deleted{cette} \deleted{\unsure{du b) et}} Nous dirons que la situation \deleted{du} multigraphe $M$ est \uline{homogène}, s'il satisfait (59) \inserted{et si \uuline{de plus}} \deleted{les conditions indiquées, de préférence si $R\ne\varnothing$,} \deleted{que} $G_0$ est transitif sur $S$ et sur $\check S$ \unsure{.|;} \deleted{et par suite} \deleted{suite que si} \inserted{alors} un $\lambda\in Z$ \deleted{\unsure{opère}} a un pt fixe dans $S$ (resp $\check S$), il opère trivialement dans $S$ (resp. ds $\check S$), donc il est trivial si $\Gamma$ est séparé en $\check S$ (resp. en $S$). Donc

(60) Si $\Gamma$ séparé en $\check S$ (resp $S^{\unsure{*}}$), alors sous les conditions d'homogénéité (55), $Z$ opère librement sur $S$ (resp. $\check S$).

On a aussi

(61) (55) implique (41), \deleted{\unsure{avec les}} \inserted{et ses} formes équivalentes $(41')\ (41'')\ (44)\ (48)$ --- donnant donc lieu aux interprétations (48.1) à (48.8.)

\deleted{(avec 59\unsure{.|a}) Ce qui s'écrit encore (sans référence à $\pi$, en termes seulement de $G_0,Z\subset\mathcal A$ deux sous-groupes de $\mathcal A=\mathrm{Aut}\,\Gamma$ commutant l'un à l'autre, sous la forme 41}
\[
\deleted{(\unsure{*})\qquad (s_0,\check s_0)\in R\text{ et }(s_0,\lambda\check s_0)\in R,\quad\lambda\in Z\quad\Longrightarrow\lambda=1}
\]
\deleted{en effet il existe $g\in G_0$ tel que $g(s_0,\check s_0)=(s_0,\lambda\check s_0)$,}

% archive sheet number: 238
